% Figure peps_dyadic_block_guide.
% Owner: chapter/ch24_peps_dyadic_geometry.tex (def:peps_dyadic_block_guide).
% Used by: gapped-2d-notes PP Sec. 8.
% Source: OpenAI2026PolynomialPEPS, sections/06-geometry.tex:60--100 (the
% passage from 2n-squares to n-squares); no source figure.  The order in which
% the blocks are repainted is not fixed by the source; this is one stage.
% Ink-to-index: each box is one n-block of the root square [0,N]^2 with N=4n,
% carrying its block label inscribed; the four gray contours are the
% 2n-parents and the black square outline is the root square.  Blue boxes
% B_1,...,B_4 are repainted blocks with their new labels; the orange box,
% labelled S beside it, is the block being repainted, still carrying its
% parent label A_2; every other block of the root square carries its parent
% label A_2, A_3 or A_4.  The asterisks, centred on their blocks, are the
% blocks outside the root square, which carry the placeholder.  The labels
% of S's neighbours across its four edges are its exterior labels C_e: B_2
% (west), A_2 (east and south) and the placeholder (north).  Boxes are
% blocks, not tensors: no wires, no contractions, no indices; the root
% outline is a closed string, not an index.
% Boundary signature: empty.
% Not to scale: the blocks are drawn as separated boxes, not abutting tiles.
% Label macros: none.
\begin{tenkzequation}
    \tenkzkernel
    \begin{tenkz}[lattice={6x6}, bonds=none, size=m] % chktex 29
        \tndeclare{species}{repainted}{hue=source:cyan}
        \tndeclare{species}{current}{hue=source:orange}
        % A free label stands half a cell north of its address, so each
        % placeholder is addressed half a cell south to centre it on its block.
        \tnmark[form=label]{0.5 s of (1,1)}{$\ast$}
        \tnmark[form=label]{0.5 s of (1,2)}{$\ast$}
        \tnmark[form=label]{0.5 s of (1,3)}{$\ast$}
        \tnmark[form=label]{0.5 s of (1,4)}{$\ast$}
        \tnmark[form=label]{0.5 s of (1,5)}{$\ast$}
        \tnmark[form=label]{0.5 s of (1,6)}{$\ast$}
        \tnmark[form=label]{0.5 s of (2,1)}{$\ast$}
        \tn[at=(2,2), skin=box, species=repainted]{B_1}
        \tn[at=(2,3), skin=box, species=repainted]{B_2}
        \tn[at=(2,4), name=dyadicCurrent, skin=box, species=current]{A_2}
        \tn[at=(2,5), skin=box]{A_2}
        \tnmark[form=label]{0.5 s of (2,6)}{$\ast$}
        \tnmark[form=label]{0.5 s of (3,1)}{$\ast$}
        \tn[at=(3,2), skin=box, species=repainted]{B_3}
        \tn[at=(3,3), skin=box, species=repainted]{B_4}
        \tn[at=(3,4), skin=box]{A_2}
        \tn[at=(3,5), skin=box]{A_2}
        \tnmark[form=label]{0.5 s of (3,6)}{$\ast$}
        \tnmark[form=label]{0.5 s of (4,1)}{$\ast$}
        \tn[at=(4,2), skin=box]{A_3}
        \tn[at=(4,3), skin=box]{A_3}
        \tn[at=(4,4), skin=box]{A_4}
        \tn[at=(4,5), skin=box]{A_4}
        \tnmark[form=label]{0.5 s of (4,6)}{$\ast$}
        \tnmark[form=label]{0.5 s of (5,1)}{$\ast$}
        \tn[at=(5,2), skin=box]{A_3}
        \tn[at=(5,3), skin=box]{A_3}
        \tn[at=(5,4), skin=box]{A_4}
        \tn[at=(5,5), skin=box]{A_4}
        \tnmark[form=label]{0.5 s of (5,6)}{$\ast$}
        \tnmark[form=label]{0.5 s of (6,1)}{$\ast$}
        \tnmark[form=label]{0.5 s of (6,2)}{$\ast$}
        \tnmark[form=label]{0.5 s of (6,3)}{$\ast$}
        \tnmark[form=label]{0.5 s of (6,4)}{$\ast$}
        \tnmark[form=label]{0.5 s of (6,5)}{$\ast$}
        \tnmark[form=label]{0.5 s of (6,6)}{$\ast$}
        \tnmark[form=enclosure]{(2,2) .. (3,3)}{}
        \tnmark[form=enclosure]{(2,4) .. (3,5)}{}
        \tnmark[form=enclosure]{(4,2) .. (5,3)}{}
        \tnmark[form=enclosure]{(4,4) .. (5,5)}{}
        \tnwire[kind=string, closed, route=orth, name=bgRoot,
            via={0.7 n of 0.7 w of (2,2), 0.7 n of 0.7 e of (2,5),
                    0.7 s of 0.7 e of (5,5), 0.7 s of 0.7 w of (5,2)}]
        \tnmark[form=label, species=current, label pos=e]{(2,4)}{$S$}
    \end{tenkz}
\end{tenkzequation}
